\clearpage\begingroup\SourceLeaf{391b}\ArithmeticHead{60}{L}{STÄCKEL, GAUSS ALS GEOMETER.}
\renewcommand{\thefootnote}{\arabic{footnote})}\setcounter{footnote}{0}
\begin{center}
\textbf{15.}\par\smallskip
\textls[100]{Elliptische, im besonderen lemniskatische Funktionen.}
\end{center}
\noindent\hspace*{\parindent}Ein zweiter Anlass, sich mit den imaginären Grössen zu beschäftigen,
eröffnete sich für \textsc{Gauss} in der doppelten Periodizität der lemniskatischen
Funktionen. Im Januar 1797 hat er diese Funktionen zu betrachten begonnen
(T. Nr. 51) und ist spätestens im März zur Entdeckung der zweiten,
imaginären Periode gelangt. Somit ergab sich »die Notwendigkeit, das Gebiet
einer veränderlichen Grösse dadurch zu erweitern, dass dieser Grösse auch
komplexe Werte beigelegt werden« (\textsc{Schlesinger}, S. 12). Die darin liegenden
Schwierigkeiten kamen sogleich zum Vorschein, als \textsc{Gauss}, die lemniskatischen
Funktionen mit dem arithmetisch-geometrischen Mittel verknüpfend, Ende
1797 zu dem allgemeinen elliptischen Integral erster Gattung überging. Die
Realitätsverhältnisse der Perioden sind ihm erst allmählich klar geworden.
Bezeichnend hierfür ist eine Aufzeichnung, die, wie es scheint, aus dem Anfang
des Jahres 1800 stammt: »Der Radikalfehler, woran meine bisherigen
Bestrebungen, den Geist der elliptischen Funktion zu verkörpern, gescheitert
sind, scheint der zu sein, dass ich dem Integral
\[
\int\frac{d\varphi}{\surd(1-e^2\sin\varphi^2)}
\]
\noindent die Bedeutung als Ausdruck eines endlichen Teils der Kugelfläche habe unterlegen
wollen, während es wahrscheinlich nur einen unendlich schmalen
\clearpage\SourceLeaf{392}\ArithmeticHead{61}{R}{DIE KOMPLEXEN GRÖSSEN IN IHRER BEZIEHUNG ZUR GEOMETRIE.}\setcounter{footnote}{0}
\noindent Kugelsektor ausdrückt« (W. X\textsubscript{1}, S. 546). Offenbar bedeutet, wie \textsc{Schlesinger} dazu
bemerkt, Kugelfläche den Ort der komplexen Veränderlichen, der endliche
Teil, dessen Ausdruck das Integral sein sollte, das Bild des Periodenparallelogramms,
während man zu einem unendlich schmalen Kugelsektor gelangen
würde, wenn das Verhältnis der Perioden reell ausfiele; vgl. Werke X\textsubscript{1}, S. 515.
\par
In das Jahr 1800 fallen auch Untersuchungen über das arithmetisch-geometrische
Mittel (T. Nr. 109). \textsc{Gauss} hat damals die wesentlichen Eigenschaften
der elliptischen Modulfunktion aufgefunden; das aber war nur möglich,
wenn er den Bereich der Veränderlichen auf das komplexe Gebiet ausdehnte.
Man wird daher behaupten dürfen, dass die Auffassungen, die er in
dem Briefe an \textsc{Bessel} vom 18. Dezember 1811 (W. VIII, S. 90, X\textsubscript{1}, S. 366)
ausgesprochen hat, bis in die Zeit um 1800 zurückreichen. Er verlangt hier,
dass man bei der Einführung einer neuen Funktion in die Analysis erkläre,
ob man sich auf reelle Werte der Veränderlichen beschränke oder seinem
Grundsatze beitrete, »dass man in dem Reiche der Grössen die imaginären
\(a+ib\) als gleiche Rechte mit den reellen geniessend ansehen müsse«. Es
folgen Auseinandersetzungen über den Sinn des Integrals bei Funktionen
einer komplexen Veränderlichen, Dabei, sagt \textsc{Gauss}, man könne »das \textls[100]{ganze}
Reich \textls[100]{aller} Grössen, reeller und imaginärer Grössen, sich durch eine unendliche
Ebene sinnlich machen, worin jeder Punkt, durch Abszisse \(=a\), Ordinate
\(=b\) bestimmt, die Grösse \(a+ib\) gleichsam repräsentiert«. Dies ist die
erste uns bekannte Stelle, wo er die geometrische Versinnlichung der komplexen
Grössen schriftlich festgelegt hat.
\par
In dem Entwurf einer Abhandlung über die Konvergenz der Reihen, der
aus der Zeit um das Jahr 1851 stammt, hat \textsc{Gauss} seine Ansichten folgendermassen
zusammengefasst: »Die vollständige Erkenntnis der Natur einer analytischen
Funktion muss auch die Einsicht in ihr Verhalten bei den imaginären
Werten des Arguments in sich schliessen, und oft ist sogar letztere unentbehrlich
zu einer richtigen Beurteilung der Gebahrung der Funktionen im
Gebiete der reellen Argumente. Unerlässlich ist es daher auch, dass die ursprüngliche
Festsetzung des Begriffes der Funktion sich mit gleicher Bündigkeit
über das \textls[100]{ganze} Grössengebiet erstrecke, welches die reellen und die
imaginären Grössen unter dem gemeinschaftlichen Namen der komplexen
Grössen in sich begreift« (W. X\textsubscript{1}, S. 405).
\clearpage\SourceLeaf{393a}\ArithmeticHead{62}{L}{STÄCKEL, GAUSS ALS GEOMETER.}\setcounter{footnote}{0}
\noindent\hspace*{\parindent}In einem zweiten Entwurfe hat \textsc{Gauss} seine Ansichten genauer darzulegen
begonnen (W. X\textsubscript{1}, S. 407—416). Wir werden darauf in Nr. 19 eingehen und
fahren fort in der Schilderung der Frühzeit.
\clearpage\endgroup

